\section{Motivation}

\begin{frame}{From Inheritance to Capabilities}
\small
\begin{columns}[T]
\begin{column}{0.48\textwidth}
\textbf{Class inheritance}
\begin{itemize}
  \item Reuse is tied to a fixed hierarchy.
  \item Base-class changes can silently affect subclasses.
  \item Designers must decide the hierarchy early.
\end{itemize}
\end{column}
\begin{column}{0.48\textwidth}
\textbf{Protocol-oriented style}
\begin{itemize}
  \item Abstraction is based on capabilities.
  \item Conformance can be added outside the type definition.
  \item Generic code asks for required behavior, not ancestry.
\end{itemize}
\end{column}
\end{columns}

\vspace{2mm}
\begin{block}{Question}
What is the typed core meaning of Swift-style protocols, generic constraints, and \texttt{any P} values?
\end{block}
\end{frame}

\begin{frame}{Why \lang?}
\small
Swift protocols are expressive, but the full language is too large for a clean TAPL-style formalization.

\vspace{2mm}
\begin{itemize}
  \item \textbf{Surface language:} a small Swift-like fragment with structs, protocols, extensions, generic constraints, first-class polymorphic functions, and \Any{P}.
  \item \textbf{Core language:} System F with records, existentials, conditionals, and natural-number primitives.
  \item \textbf{Main idea:} protocol features are not primitive at runtime; they elaborate into ordinary typed core constructs.
\end{itemize}

\vspace{1mm}
\[
  \text{well-typed surface program}
  \quad\Longrightarrow\quad
  \text{well-typed elaborated core program}
\]
\end{frame}

\begin{frame}{Scope}
\small
\begin{center}
\begin{tabular}{lll}
\toprule
Feature area & In scope & Out of scope \\
\midrule
Data & structs, fields & classes, inheritance \\
Protocols & requirements, extensions & associated types \\
Polymorphism & parametric, constrained, first-class & implicit inference \\
Existentials & \Any{P} packages & \texttt{some P} opaque types \\
Effects & pure expressions & references, mutation, effects \\
\bottomrule
\end{tabular}
\end{center}

\vspace{2mm}
\begin{itemize}
  \item The fragment is intentionally small enough to elaborate and prove.
  \item It is still rich enough to model delegate-style and protocol-oriented Swift patterns.
\end{itemize}
\end{frame}

\section{Surface Language Design}

\begin{frame}[fragile]{Surface Syntax: Protocols and Conformance}
\small
\begin{lstlisting}
protocol Named {
  func name() -> String
}

struct Screen {
  let identifier: Int
}

extension Screen: Named {
  func name() -> String { "Home" }
}
\end{lstlisting}

\vspace{1mm}
\begin{itemize}
  \item Protocols declare required methods.
  \item Extensions provide conformance evidence for a concrete struct.
  \item Methods are normalized internally by inserting an explicit \texttt{self} parameter.
\end{itemize}
\end{frame}

\begin{frame}{Surface Fragment}
\small
\[
\begin{array}{rcll}
\tau &::=& \Int \mid \Bool \mid \String \mid \Unit
  & \text{base types} \\
&\mid& S \mid X \mid \Self
  & \text{structs, type variables, contextual self} \\
&\mid& \Any{P}
  & \text{existential protocol value} \\
&\mid& \ftype{\overline{X:\overline{P}}}{\overline{\ell:\tau}}{\tau}
  & \text{function scheme}
\end{array}
\]

\vspace{1mm}
\[
\begin{array}{rcl}
d &::=&
\kw{struct}\ S\ \{\overline{\kw{let}\ x:\tau};\overline{md}\}
\\
&\mid&
\kw{protocol}\ P\ \{\overline{m:\sigma}\}
\\
&\mid&
\kw{extension}\ S:P\ \{\overline{md}\}
\\
&\mid&
\kw{func}\ f\langle\overline{X:\overline{P}}\rangle
(\overline{\ell\ x:\tau})\to\tau\ \{ss\}
\end{array}
\]
\end{frame}

\begin{frame}{Typing Environments}
\small
\[
\Psi=(\Psi_S,\Psi_P,\Psi_C,\Psi_M,\Psi_V)
\]

\begin{center}
\begin{tabular}{ll}
\toprule
Component & Meaning \\
\midrule
$\Psi_S(S)$ & fields of struct $S$ \\
$\Psi_P(P)$ & method requirements of protocol $P$ \\
$\Psi_C(S,P)$ & conformance methods for $S:P$ \\
$\Psi_M(S)$ & methods available on struct $S$ \\
$\Psi_V(x)$ & top-level values and functions \\
\bottomrule
\end{tabular}
\end{center}

\vspace{1mm}
\begin{itemize}
  \item $\Omega$ records generic constraints such as $X:P$.
  \item $\Gamma$ records local values.
  \item $\rho$ resolves \Self\ in struct, extension, and protocol contexts.
\end{itemize}
\end{frame}

\begin{frame}{Protocol Satisfaction}
\small
Only two surface types satisfy a protocol:

\[
\inferrule*[Right=\textsc{Sat-Struct}]
  {(S,P)\in\dom(\Psi_C)}
  {\Psi;\Omega\vdash S\modelsP P}
\qquad
\inferrule*[Right=\textsc{Sat-TVar}]
  {X:\overline{P}\in\Omega \\ P\in\overline{P}}
  {\Psi;\Omega\vdash X\modelsP P}
\]

\vspace{2mm}
\begin{itemize}
  \item Concrete structs satisfy protocols through extensions.
  \item Type variables satisfy protocols through generic constraints.
  \item \Any{P} does \emph{not} itself satisfy $P$; it hides a witness type and its dictionary.
\end{itemize}
\end{frame}

\section{Elaboration Rules}

\begin{frame}{Core Target}
\small
\begin{block}{Target calculus}
System F + records + existentials + conditionals + recursion and natural-number primitives.
\end{block}

\[
\begin{array}{rcl}
A &::=& \Nat \mid \Bool \mid \String \mid \Unit
       \mid X \mid A\to A \mid \forall X.A
       \mid \{\overline{l:A}\} \mid \exists X.A
\\[1mm]
t &::=& x \mid \lambda x:A.t \mid t\ t
       \mid \Lambda X.t \mid t[A]
       \mid \{\overline{l=t}\} \mid t.l
       \mid \kw{pack} \mid \kw{unpack}\ \ldots
\end{array}
\]

\vspace{1mm}
Elaboration removes the protocol layer by making all evidence explicit.
\end{frame}

\begin{frame}{Type Translation}
\small
\[
\begin{array}{rcl}
\trans{\Int} &=& \Nat \\
\trans{\Bool} &=& \Bool \\
\trans{\String} &=& \String \\
\trans{S} &=& \{\overline{x:\trans{\tau}}\}
  \quad \text{where }\Psi_S(S)=\{\overline{x:\tau}\} \\
\trans{\Any{P}} &=&
  \exists X.\{\kw{value}:X,\kw{dict}:\Dict(P,X)\}
\end{array}
\]

\vspace{1mm}
\[
  \Dict(P,A)=\{\overline{m_i:\transSelf{A}{\sigma_i}}\}
  \quad\text{where }\Psi_P(P)=\{\overline{m_i:\sigma_i}\}.
\]

\vspace{1mm}
\begin{itemize}
  \item A protocol dictionary is a record of implementations.
  \item Abstract protocol \Self is interpreted as the receiver type $A$ inside the dictionary.
\end{itemize}
\end{frame}

\begin{frame}{Evidence Elaboration Rules}
\scriptsize
\[
  \Psi;\Omega;\Delta \has \tau \modelsP P \elab d
\]

\begin{mathpar}
\inferrule*[Right=\textsc{Ev-Struct}]
  { (S,P)\in\dom(\Psi_C) }
  { \Psi;\Omega;\Delta \has S \modelsP P \elab \delta(S,P) }

\inferrule*[Right=\textsc{Ev-TVar}]
  { X:\overline{P}\in\Omega \\ P\in\overline{P} \\ \Delta(X,P)=d }
  { \Psi;\Omega;\Delta \has X \modelsP P \elab d }
\end{mathpar}

\vspace{1mm}
\begin{itemize}
  \item Concrete conformance becomes a top-level dictionary name.
  \item Generic conformance becomes a dictionary variable from $\Delta$.
\end{itemize}
\end{frame}

\begin{frame}{Generic Functions and Calls}
\scriptsize
\begin{mathpar}
\inferrule*[Right=\textsc{E-Fun}]
  {
    \Omega'=\Omega,\overline{X_i:\overline{P_i}} \\
    \Delta'=\Delta,\overline{(X_i,P)=\delta(X_i,P)} \\
    \Gamma_f=\Gamma,\overline{x_j:\tau_j} \\
    \Psi;\Omega';\Gamma_f;\rho;\Delta'
      \has ss:\tau_r \elab t_f
  }
  {
    \ectx \has
      \kw{func}\langle\overline{X_i:\overline{P_i}}\rangle
      (\overline{\ell_j\ x_j:\tau_j})\to\tau_r\ \{ss\}
      :
      \ftype{\overline{X_i:\overline{P_i}}}
             {\overline{\ell_j:\tau_j}}
             {\tau_r}
      \elab M_{\lambda}
  }
\end{mathpar}

\[
M_{\lambda} =
  \Lambda \overline{X_i}.\
  \lambda \overline{\delta(X_i,P):\Dict(P,X_i)}.\
  \lambdastar\ \overline{x_j:\trans{\tau_j}}.\ t_f
\]

\begin{mathpar}
\inferrule*[Right=\textsc{E-TypeApp}]
  {
    \ectx \has e:
      \ftype{\overline{X_i:\overline{P_i}}}
             {\overline{\ell_j:\tau_j}}
             {\tau_r}
      \elab t \\
    \overline{\Psi;\Omega;\Delta \has \tau_i' \modelsP P
      \elab d_{i,P}}
  }
  {
    \ectx \has e\langle\overline{\tau_i'}\rangle:
      \subst{\overline{\tau_i'}}{\overline{X_i}}
      \left(
        \ftype{\emptyctx}{\overline{\ell_j:\tau_j}}{\tau_r}
      \right)
      \elab
      t[\overline{\trans{\tau_i'}}]\ \overline{d_{i,P}}
  }
\end{mathpar}
\end{frame}

\begin{frame}{Method Elaboration Rules}
\scriptsize
\begin{mathpar}
\inferrule*[Right=\textsc{E-StructMethod}]
  {
    \ectx \has e:S \elab t \\
    \Psi_M(S)(m)=
      \ftype{\overline{X_i:\overline{P_i}}}
             {\epsilon:S,\overline{\ell_j:\tau_j}}
             {\tau_r}
  }
  {
    \ectx \has e.m:
      \ftype{\overline{X_i:\overline{P_i}}}
             {\overline{\ell_j:\tau_j}}
             {\tau_r}
      \elab M_S
  }

\inferrule*[Right=\textsc{E-ConstrainedMethod}]
  {
    \ectx \has e:X \elab t \\
    X:\overline{P}\in\Omega \\
    P_k\in\overline{P} \\
    \Psi_P(P_k)(m)=
      \ftype{\overline{Y_i:\overline{Q_i}}}
             {\epsilon:\Self,\overline{\ell_j:\tau_j}}
             {\tau_r} \\
    \Delta(X,P_k)=d
  }
  {
    \ectx \has e.m:
      [\Self\mapsto X]
      \left(
        \ftype{\overline{Y_i:\overline{Q_i}}}
               {\overline{\ell_j:\tau_j}}
               {\tau_r}
      \right)
      \elab M_X
  }
\end{mathpar}

\[
\begin{aligned}
M_S &=
  \Lambda \overline{X_i}.\
  \lambda \overline{\delta(X_i,P):\Dict(P,X_i)}.\
  \lambdastar\ \overline{y_j:\trans{\tau_j}}.\
  \appstar(
    \mu(S,m)[\overline{X_i}]\ \overline{\delta(X_i,P)}\ t,
    \overline{y_j}),\\
M_X &=
  \Lambda \overline{Y_i}.\
  \lambda \overline{\delta(Y_i,Q):\Dict(Q,Y_i)}.\
  \lambdastar\ \overline{y_j:\trans{[\Self\mapsto X]\tau_j}}.\
  \appstar(
    d.m[\overline{Y_i}]\ \overline{\delta(Y_i,Q)}\ t,
    \overline{y_j})
\end{aligned}
\]
\end{frame}

\begin{frame}{Existential Elaboration Rules}
\tiny
\begin{mathpar}
\inferrule*[Right=\textsc{E-AsAny}]
  {
    \ectx \has e:\tau \elab t \\
    \Psi;\Omega;\Delta \has \tau \modelsP P \elab d
  }
  {
    \ectx \has e\ \kw{as}\ \Any{P}:\Any{P}
      \elab
      \kw{pack}[\trans{\tau},\{\kw{value}=t,\kw{dict}=d\}]
      \ \kw{as}\
      \exists X.\{\kw{value}:X,\kw{dict}:\Dict(P,X)\}
  }
\end{mathpar}

\begin{mathpar}
\inferrule*[Right=\textsc{E-ExistentialMethod}]
  {
    \ectx \has e:\Any{P}\elab t \\
    \Psi_P(P)(m)=
      \ftype{\overline{Y_i:\overline{Q_i}}}
             {\epsilon:\Self,\overline{\ell_j:\tau_j}}
             {\tau_r} \\
    \Self\notin\fv(\overline{\tau_j}) \\
    \Self\notin\fv(\tau_r)
  }
  {
    \ectx \has e.m:
      \ftype{\overline{Y_i:\overline{Q_i}}}
             {\overline{\ell_j:\tau_j}}
             {\tau_r}
      \elab M_{\exists}
  }
\end{mathpar}

\[
\begin{gathered}
M_{\exists} =
  \Lambda \overline{Y_i}.\
  \lambda \overline{\delta(Y_i,Q):\Dict(Q,Y_i)}.\
  \lambdastar\ \overline{y_j:\trans{\tau_j}}.\\
  \kw{unpack}\ \{X,z\}=t\ \kw{in}\
  \appstar(
    z.\kw{dict}.m[\overline{Y_i}]\
      \overline{\delta(Y_i,Q)}\
      z.\kw{value},
    \overline{y_j})
\end{gathered}
\]
\end{frame}

\section{Examples}

\begin{frame}[fragile]{Example 1: Constrained Generic}
\small
\begin{lstlisting}
protocol Named {
  func name() -> String
}

struct Screen {
  let identifier: Int
}

extension Screen: Named {
  func name() -> String { "Home" }
}

func keepNamed<T: Named>(item: T) -> T {
  item
}

let screen: Screen = Screen(identifier: 1)
let kept: Screen = keepNamed<Screen>(item: screen)
kept.name()
\end{lstlisting}
\end{frame}

\begin{frame}{Elaboration Shape for Example 1}
\small
\[
\begin{array}{rcl}
\Dict(\kw{Named},A)
&=&
\{\kw{name}: A\to\String\}
\\[1mm]
\kw{dict\_Screen\_Named}
&=&
\{\kw{name}=\lambda\kw{self}:\trans{\kw{Screen}}.\ "Home"\}
\\[1mm]
\kw{keepNamed}
&=&
\Lambda T.\lambda d_{T,\kw{Named}}:\Dict(\kw{Named},T).
   \lambda\kw{item}:T.\ \kw{item}
\end{array}
\]

\vspace{2mm}
\[
\begin{aligned}
&\kw{keepNamed}\langle\kw{Screen}\rangle(\kw{item}:\kw{screen})
\\
&\quad\elab\quad
\kw{keepNamed}\ [\trans{\kw{Screen}}]\ \kw{dict\_Screen\_Named}\ \kw{screen}
\end{aligned}
\]
\end{frame}

\begin{frame}[fragile]{Example 2: Existential Delegate}
\small
\begin{lstlisting}
protocol SelectionDelegate {
  func didSelect(row: Int) -> Bool
}

struct EvenSelection {
  let enabled: Bool
}

extension EvenSelection: SelectionDelegate {
  func didSelect(row: Int) -> Bool {
    if self.enabled { row < 2 } else { false }
  }
}

let concrete: EvenSelection = EvenSelection(enabled: true)
let delegate: any SelectionDelegate =
  concrete as any SelectionDelegate

delegate.didSelect(row: 0)
\end{lstlisting}
\end{frame}

\begin{frame}{Elaboration Shape for Example 2}
\small
\[
\begin{aligned}
\trans{\Any{\kw{SelectionDelegate}}}
=
\exists X.\{&
\kw{value}:X,\\
&\kw{dict}:\{\kw{didSelect}:X\to\Nat\to\Bool\}\}
\end{aligned}
\]

\vspace{1mm}
The cast packages the concrete object with its conformance evidence:
\[
\begin{aligned}
\kw{concrete}\ \kw{as}\ \Any{\kw{SelectionDelegate}}
&\elab
\kw{pack}\bigl[\\
\trans{\kw{EvenSelection}},
&\{\kw{value}=\kw{concrete},\\
&\ \kw{dict}=\kw{dict\_EvenSelection\_SelectionDelegate}\}\bigr]
\end{aligned}
\]

\vspace{1mm}
The method call opens the package and calls the dictionary method:
\[
\begin{aligned}
&\kw{unpack}\ \{X,z\}=\kw{delegate}\ \kw{in}\\
&\quad z.\kw{dict}.\kw{didSelect}\ z.\kw{value}\ 0
\end{aligned}
\]
\end{frame}

\begin{frame}[fragile]{Example 3: Settings Screen}
\small
\begin{lstlisting}
struct SettingsViewController {
  let titleProvider: any ScreenTitleProvider
  let dataSource: any TableViewDataSource
  let delegate: any TableViewDelegate
  let reuse: ReuseIdentity

  func screenTitle() -> String {
    self.titleProvider.title()
  }

  func canOpenFirstCell() -> Bool {
    self.delegate.didSelect(row: 0)
  }

  func stableIndex() -> Int {
    self.reuse.keep<Int>(0)
  }
}
\end{lstlisting}

\vspace{1mm}
This combines existential delegates, constrained protocol calls, and a first-class polymorphic function stored in a struct field.
\end{frame}

\section{Elaboration Safety}

\begin{frame}{Safety Statement}
\small
\begin{block}{Elaboration preservation}
If a surface expression, statement sequence, declaration, or program type checks and elaborates to core, then the generated core term has the translated type.
\end{block}

\[
\Psi;\Omega;\Gamma;\rho\vdash e:\tau
\qquad
\Psi;\Omega;\Gamma;\rho;\Delta\vdash e:\tau\elab t
\]
\[
\Longrightarrow
\qquad
\mathcal{E}\vdash_c t:\trans{\tau}.
\]

\vspace{1mm}
\[
  \mathcal{E}
  =
  \mathcal{G}_\Psi;\trans{\Omega};\trans{\Gamma};\trans{\Delta}
\]

\vspace{2mm}
\begin{itemize}
  \item $\mathcal{G}_\Psi$ assigns translated types to top-level values, method functions, and conformance dictionaries.
  \item $\trans{\Delta}$ contains the dictionary variables required by active generic constraints.
  \item Once the generated core program is well typed, safety follows from standard core progress and preservation.
\end{itemize}
\end{frame}

\begin{frame}{Proof Roadmap}
\small
\begin{columns}[T]
\begin{column}{0.38\textwidth}
\centering
\[
\begin{array}{c}
\text{evidence preservation}
\\[0.5mm]\Downarrow\\[-1mm]
\text{expression preservation}
\\[0.5mm]\Downarrow\\[-1mm]
\text{statement-sequence preservation}
\\[0.5mm]\Downarrow\\[-1mm]
\text{declaration preservation}
\\[0.5mm]\Downarrow\\[-1mm]
\text{program preservation}
\\[0.5mm]\Downarrow\\[-1mm]
\text{end-to-end safety}
\end{array}
\]
\end{column}

\begin{column}{0.58\textwidth}
\scriptsize
\begin{center}
\begin{tabular}{ll}
\toprule
Lemma family & Role \\
\midrule
Type/context & translated contexts are well formed \\
Substitution & type application commutes with translation \\
\Self substitution & methods match dictionary fields \\
Arguments & labels keep parameter types \\
Wrappers & dictionaries become parameters \\
\bottomrule
\end{tabular}
\end{center}
\end{column}
\end{columns}

\vspace{2mm}
\begin{itemize}
  \item Each surface feature is reduced to ordinary core typing rules.
\end{itemize}
\end{frame}

\begin{frame}{Evidence and Declarations}
\small
\begin{block}{Theorem 1: evidence preservation}
If surface evidence proves $\tau \modelsP P$ and elaborates to $d$, then
\[
  \mathcal{G}_\Psi;\trans{\Omega};\trans{\Delta}
  \vdash_c d:\Dict(P,\trans{\tau}).
\]
\end{block}

\begin{itemize}
  \item Concrete evidence uses the global dictionary
  $\delta(S,P):\Dict(P,\trans{S})$.
  \item Generic evidence uses the local dictionary variable
  $\Delta(X,P)=d$, so $d:\Dict(P,X)$.
  \item Declaration preservation proves extension dictionaries:
  \[
    \delta(S,P)=\{\overline{m_i=\mu(S,m_i)}\}.
  \]
  \item Lemma 5 turns $[\Self\mapsto S]\sigma_i$ into the field type expected by $\Dict(P,\trans{S})$.
\end{itemize}
\end{frame}

\begin{frame}{Expression Preservation: Main Cases}
\small
\begin{center}
\begin{tabular}{p{0.35\textwidth}p{0.55\textwidth}}
\toprule
Case & Proof idea \\
\midrule
Generic application
&
Type-apply the callee, then pass dictionaries supplied by Theorem 1. Lemma 4 rewrites the result type after substitution.
\\[1mm]
Function expression
&
Turn constraints into dictionary parameters before ordinary parameters. Lemma 8 types the generated wrapper.
\\[1mm]
Concrete method
&
Use global method function $\mu(S,m)$, apply captured receiver, then visible arguments.
\\[1mm]
Constrained method
&
Project from evidence dictionary $d.m$; Lemma 5 identifies protocol \Self with the receiver type variable.
\\[1mm]
Constructor/field
&
Structs translate to records, so introduction and projection are direct core typing rules.
\\
\bottomrule
\end{tabular}
\end{center}
\end{frame}

\begin{frame}{Existential Method Case}
\small
\[
\trans{\Any{P}}
=
\exists X.\{\kw{value}:X,\kw{dict}:\Dict(P,X)\}
\]

\vspace{1mm}
For a call $e.m(\overline{a})$, the generated core term has the shape
\[
\kw{unpack}\ \{X,z\}=t\ \kw{in}\
z.\kw{dict}.m\ z.\kw{value}\ \overline{u}.
\]

\vspace{1mm}
\begin{itemize}
  \item Opening the package reveals a fresh hidden type $X$ only inside the unpack body.
  \item The dictionary method is typed with protocol \Self interpreted as $X$.
  \item The unpack body must have a result type that does not mention $X$.
\end{itemize}

\vspace{1mm}
A method such as
\[
\kw{clone} : \Self \to \Self
\]
would return the hidden witness type $X$ after unpacking.

\vspace{1mm}
\begin{itemize}
  \item The surface type has no name for that $X$.
  \item Lemma 6 fixes the safe case: if visible parameter and result types are \Self-free, then $\transSelf{X}{\tau}=\trans{\tau}$.
\end{itemize}
\end{frame}

\section{Implementation Details}

\begin{frame}{Implementation Pipeline}
\small
\[
\begin{array}{c}
\texttt{.swift-like source}
\to
\text{lexer/parser}
\to
\text{resolver}
\to
\text{checker}
\\[1mm]
\to
\text{elaborator}
\to
\text{core parser}
\to
\text{core typechecker/evaluator}
\end{array}
\]

\vspace{2mm}
\begin{itemize}
  \item Implemented in MoonBit.
  \item Surface parser produces an AST close to the user syntax.
  \item Elaboration currently renders readable core concrete syntax, then reuses the existing core parser and interpreter.
\end{itemize}
\end{frame}

\begin{frame}{Three Surface Passes}
\small
\begin{center}
\begin{tabular}{lll}
\toprule
Pass & Input/output & Main responsibility \\
\midrule
Resolver & syntax AST $\to$ resolved AST & names, \Self, method receivers \\
Checker & resolved AST $\to$ typed AST & conformance, calls, existentials \\
Elaborator & typed AST $\to$ core syntax & dictionaries, packages, unpacking \\
\bottomrule
\end{tabular}
\end{center}

\vspace{2mm}
\begin{itemize}
  \item The checker records call and member resolution, so elaboration does not guess.
  \item The core interpreter then independently type checks the generated program.
\end{itemize}
\end{frame}

\begin{frame}{Current Limitations}
\small
\begin{itemize}
  \item Single-file programs only.
  \item Generic instantiation is explicit; no type argument inference.
  \item \Int elaborates to core \Nat, so subtraction is saturating.
  \item No string concatenation, because the current core has no string concat primitive.
  \item Existential elimination is surface-restricted to safe method calls, not explicit user-level unpack.
\end{itemize}

\vspace{2mm}
\begin{block}{Run examples}
\texttt{moon run src -- examples/<file>.swift}
\end{block}
\end{frame}

\section{Conclusion}

\begin{frame}{Takeaways}
\small
\begin{itemize}
  \item \lang\ isolates a proof-relevant subset of Swift-style protocol-oriented programming.
  \item Protocol conformance becomes explicit dictionary evidence.
  \item \Any{P} becomes a standard existential package.
  \item The preservation proof connects surface typing to core typing.
  \item The implementation follows the proof structure: resolve, check, elaborate, then run core.
\end{itemize}
\end{frame}

\begin{frame}[plain]
\centering
\vspace{20mm}
{\Large Thank you!}

\vspace{8mm}
{\small Questions?}
\end{frame}
